% GENERATED FILE. Edit canonical lessons, then run npm run generate:books.
% canonical-source-hash: fnv1a64:23c29fecbf034171
% canonical-lessons: ZH-C60-bai, ZH-W60-bai-white, ZH-C60-yanse, ZH-C60-hong, ZH-C60-hei, ZH-C60-lan

\chapter{Colours}
\label{ch:zh-colours}
\hlchaptermodality{60}

\begin{chapteropening}
\textbf{By the end of this chapter:} \emph{I can name white, red, black and blue, and ask about colours.}

Complete the last lesson of chapter 60: I can name white, red, black and blue, and ask about colours.
\end{chapteropening}

\section[\texorpdfstring{\zh{白} (bái)}{白 (bái)}]{\zh{白} (bái) — white; pale}

\label{lesson:ZH-C60-bai}



\begin{quote}
Before the new one: say the Mandarin for trousers, then the Mandarin for pretty.
\end{quote}

\subsection*{You'll want to know: \zh{白}}
\textbf{\zh{白}} — \emph{bái} — \textquotedblleft{}white; pale\textquotedblright{}.

\begin{quote}
\textbf{\zh{他的手很白}} \emph{tā de shǒu hěn bái} — His hands are very pale.
\end{quote}

The new character, in pieces: \zh{白} — a short fall on top of \zh{日}. The next lesson puts it in your hand.

\begin{cousinweb}
\textbf{\zh{白}} (\emph{bái}) is inherited from Middle Chinese \emph{baek} (Guangyun, in Baxter's transcription), glossed 'the colour of the west', the direction that the old scheme of five colours paired with white.
\end{cousinweb}

\begin{grammarlens}[title={What you've built}]
One more word for everyday answers.
\end{grammarlens}

\subsection*{Your turn}
\begin{itemize}
\raggedright
  \item \emph{Say it:} \emph{bái}
  \item \emph{Say it:} \emph{bái}, once more
  \item \emph{Say it:} \emph{tā de shǒu hěn bái} — His hands are very pale.
  \item \emph{Recall:} say the Mandarin for trousers, then the Mandarin for pretty, then say \emph{bái} again
\end{itemize}

\begin{tcolorbox}[breakable,colback=teal!4,colframe=teal!35!black,title={Before you move on}]
What does \zh{白} mean? (\textquotedblleft{}White; pale\textquotedblright{}.) Say it once more.
\end{tcolorbox}
\section[\texorpdfstring{\zh{白} (bái)}{白 (bái)}]{\zh{白} — a short fall over \zh{日}}

\label{lesson:ZH-W60-bai-white}



\begin{quote}
The word first, then the character inside it.

\begin{itemize}
\raggedright
  \item \emph{Recall:} say \emph{bái} — \textbf{R1}, one lesson back, before it decays
\end{itemize}

\begin{quote}
\zh{日} \zh{百}
\end{quote}

\emph{Read it:} each one aloud, with nothing said first — \emph{rì}, \emph{bǎi}
\end{quote}

\subsection*{Script — a short fall over \zh{日}}
\hlblockfigure{\detokenize{figures/ZH-W60-bai-white-filmstrip.pdf}}{How \zh{白} is written, stroke by stroke}

\begin{quote}
\zh{白}
\end{quote}

\textbf{Two pieces, top to bottom.}

\textbf{A short left-falling stroke} at the top, from the right down to the left.

\textbf{\zh{日}, four strokes}, the same four your hand makes in \zh{日}: the left side, the top that turns down the right side, the middle bar, and the bar that closes the bottom.

1 + 4 = five.

\subsection*{Your turn}
\begin{itemize}
\raggedright
  \item \emph{Write it:} \zh{白}, all five, with the model in front of you
  \item \emph{Cover:} write it again from memory, and count the pieces, not the strokes
\end{itemize}

\begin{tcolorbox}[breakable,colback=teal!4,colframe=teal!35!black,title={Before you move on}]
How many strokes in \zh{白}? (\textbf{Five}.) What comes before \zh{日}? (\textbf{A short fall at the top}.) Which number you know has \zh{白} under its top bar? (\textbf{\zh{百}}.)

Source: \href{https://raw.githubusercontent.com/chanind/hanzi-writer-data/68d10a4b21150cae5e1ebbd223eed289cf32d90c/data/\%E7\%99\%BD.json}{Hanzi Writer Data \zh{白}, pinned stroke paths and medians}.
\end{tcolorbox}
\section[yánsè]{yánsè — colour}

\label{lesson:ZH-C60-yanse}



\begin{quote}
Before the new one: say the Mandarin for pretty, then the Mandarin for white.
\end{quote}

\subsection*{You'll want to know: yánsè}
\emph{yánsè} — \textquotedblleft{}colour\textquotedblright{}.

\begin{quote}
\emph{wǒ ài zhège yánsè} — I love this colour.
\end{quote}

Its characters come later; for now it is a word for the ear.

\begin{cousinweb}
\emph{yánsè} is built from yán 'face' and sè, both inherited from Middle Chinese, \emph{ngaen} and \emph{srik} (Guangyun, in Baxter's transcription). The dictionary glosses the second with the whole word, 'colour'.
\end{cousinweb}

\begin{grammarlens}[title={What you've built}]
One more word for everyday answers.
\end{grammarlens}

\subsection*{Your turn}
\begin{itemize}
\raggedright
  \item \emph{Say it:} \emph{yánsè}
  \item \emph{Say it:} \emph{yánsè}, once more
  \item \emph{Say it:} \emph{wǒ ài zhège yánsè} — I love this colour.
  \item \emph{Recall:} say the Mandarin for pretty, then the Mandarin for white, then say \emph{yánsè} again
\end{itemize}

\begin{tcolorbox}[breakable,colback=teal!4,colframe=teal!35!black,title={Before you move on}]
What does yánsè mean? (\textquotedblleft{}Colour\textquotedblright{}.) Say it once more.
\end{tcolorbox}
\section[hóng]{hóng — red}

\label{lesson:ZH-C60-hong}



\begin{quote}
Before the new one: say the Mandarin for white, then the Mandarin for colour.
\end{quote}

\subsection*{You'll want to know: hóng}
\emph{hóng} — \textquotedblleft{}red\textquotedblright{}.

\begin{quote}
\emph{zhège píngguǒ hěn hóng} — This apple is really red.
\end{quote}

Its characters come later; for now it is a word for the ear.

\begin{cousinweb}
\emph{hóng} is inherited from Middle Chinese \emph{huwng} (Guangyun, in Baxter's transcription), glossed 'a colour'; which colour is not said there.
\end{cousinweb}

\begin{grammarlens}[title={What you've built}]
One more word for everyday answers.
\end{grammarlens}

\subsection*{Your turn}
\begin{itemize}
\raggedright
  \item \emph{Say it:} \emph{hóng}
  \item \emph{Say it:} \emph{hóng}, once more
  \item \emph{Say it:} \emph{zhège píngguǒ hěn hóng} — This apple is really red.
  \item \emph{Recall:} say the Mandarin for white, then the Mandarin for colour, then say \emph{hóng} again
\end{itemize}

\begin{tcolorbox}[breakable,colback=teal!4,colframe=teal!35!black,title={Before you move on}]
What does hóng mean? (\textquotedblleft{}Red\textquotedblright{}.) Say it once more.
\end{tcolorbox}
\section[hēi]{hēi — black; dark}

\label{lesson:ZH-C60-hei}



\begin{quote}
Before the new one: say the Mandarin for colour, then the Mandarin for red.
\end{quote}

\subsection*{You'll want to know: hēi}
\emph{hēi} — \textquotedblleft{}black; dark\textquotedblright{}.

\begin{quote}
\emph{nàge fángjiān hěn hēi} — That room is dark.
\end{quote}

Its characters come later; for now it is a word for the ear.

\begin{cousinweb}
\emph{hēi} is inherited from Middle Chinese \emph{xok} (Guangyun, in Baxter's transcription), glossed 'the colour of the north', which the old scheme of five colours paired with black.
\end{cousinweb}

\begin{grammarlens}[title={What you've built}]
One more word for everyday answers.
\end{grammarlens}

\subsection*{Your turn}
\begin{itemize}
\raggedright
  \item \emph{Say it:} \emph{hēi}
  \item \emph{Say it:} \emph{hēi}, once more
  \item \emph{Say it:} \emph{nàge fángjiān hěn hēi} — That room is dark.
  \item \emph{Recall:} say the Mandarin for colour, then the Mandarin for red, then say \emph{hēi} again
\end{itemize}

\begin{tcolorbox}[breakable,colback=teal!4,colframe=teal!35!black,title={Before you move on}]
What does hēi mean? (\textquotedblleft{}Black; dark\textquotedblright{}.) Say it once more.
\end{tcolorbox}
\section[lán]{lán — blue}

\label{lesson:ZH-C60-lan}



\begin{quote}
Before the new one: say the Mandarin for red, then the Mandarin for black.
\end{quote}

\subsection*{You'll want to know: lán}
\emph{lán} — \textquotedblleft{}blue\textquotedblright{}.

\begin{quote}
\emph{tā de yǎnjing hěn lán} — Her eyes are very blue.
\end{quote}

Its characters come later; for now it is a word for the ear.

\begin{cousinweb}
\emph{lán} is inherited from Middle Chinese \emph{lam} (Guangyun, in Baxter's transcription), glossed 'a plant for dye'. The colour is not given there; blue as the colour of that dye is the likely step, but it is not settled.
\end{cousinweb}

\begin{grammarlens}[title={What you've built}]
One more word for everyday answers.
\end{grammarlens}

\subsection*{Your turn}
\begin{itemize}
\raggedright
  \item \emph{Say it:} \emph{lán}
  \item \emph{Say it:} \emph{lán}, once more
  \item \emph{Say it:} \emph{tā de yǎnjing hěn lán} — Her eyes are very blue.
  \item \emph{Recall:} say the Mandarin for red, then the Mandarin for black, then say \emph{lán} again
\end{itemize}

\begin{tcolorbox}[breakable,colback=teal!4,colframe=teal!35!black,title={Before you move on}]
What does lán mean? (\textquotedblleft{}Blue\textquotedblright{}.) Say it once more.
\end{tcolorbox}
